\scanpage{24}\mspage{9}
\uline{et pour différentielle invariante} $\pi=\mathrm{dx}/2\mathrm{y}$.

De plus, sous ces hypothèses, la donnée d'une forme différentielle invariante suffit à rigidifier une courbe.

4.\qquad \uline{Cas où} $2$ \uline{est nilpotent et} $\mathrm{c}_{4}$ \uline{inversible}.

Lorsque $2=0$, la forme différentielle ne suffit plus jamais à rigidifier une courbe de genre un. Lorsque $\mathrm{c}_{4}$ est inversible et $2$ nilpotent, donc $\mathrm{a}_{1}$ inversible, en changeant le choix de $\pi$, $\mathrm{x}$, $\mathrm{y}$ on peut trouver une équation
\[
\mathrm{y}^{2}+\mathrm{xy}=\mathrm{x}^{3}+\mathrm{a}_{2}\,\mathrm{x}^{2}+\mathrm{a}_{6}\qquad
\pi=\handwritten{\frac}{\mathrm{dx}}{\mathrm{x}+2\mathrm{y}}\qquad.
\tag*{(4.1)}
\]

Les autres choix de $\mathrm{x}$ et $\mathrm{y}$ donnant lieu à une telle équation s'obtiennent en faisant, dans (1.2)
\[
\mathrm{r}=\mathrm{t}=0\qquad\mathrm{u}=1+2\mathrm{s}
\tag*{(4.2)}
\]
et on aura
\[
\begin{aligned}[t]
\mathrm{u}^{2}\,\mathrm{a}'_{2}&=\mathrm{a}_{2}-\mathrm{s}-\mathrm{s}^{2}\\
\mathrm{u}^{6}\,\mathrm{a}'_{6}&=\mathrm{a}_{6}
\end{aligned}
\tag*{(4.3)}
\]
et \uline{modulo 2}
\[
\mathrm{c}_{4}=1\qquad\mathrm{c}_{6}=1\qquad\Delta=\mathrm{a}_{6}\qquad.
\tag*{(4.4)}
\]
\margin{\handwritten{$\unsure{\overline{\mathrm{I}}|T}\,\unsure{;|:}\,\unsure{\sim|m|n}$}}

5.\qquad \uline{Cas où} $\mathrm{c}_{4}=0$, $\mathrm{p}=2$ \uline{ou} $\mathrm{p}=3$\mbox{\begingroup\edef\encoutersep{\the\dimexpr.5\pgflinewidth\relax}\setbox0=\hbox{.}\edef\encwda{\the\wd0}\edef\enchta{\the\ht0}\edef\encdpa{\the\dp0}\rlap{\smash{\tikz[baseline=0pt]{\begin{scope}[shift={(\the\dimexpr \encwda/2\relax,\the\dimexpr -\encdpa/2+\enchta/2\relax)}]\coordinate (enc-mark-center) at (0pt,0pt);\handwritten{\draw[overlay,thin,gray] ([xshift=.1em,yshift=.86em]enc-mark-center) .. controls +(-.2em,.5em) and +(-.83em,-.8em) .. ([xshift=-.23em,yshift=-.9em]enc-mark-center) .. controls +(.67em,-.12em) and +(0,-.45em) .. ([xshift=.93em,yshift=.5em]enc-mark-center);}\end{scope}\pgfresetboundingbox\path[use as bounding box] (0pt,\the\dimexpr-\encdpa\relax) rectangle (\encwda,\enchta);}}}\copy0\endgroup}\margin{\handwritten{\tikz[baseline]{\draw[thin,gray] (0,0) .. controls (.12em,.38em) and (.38em,.27em) .. (.58em,.25em);}}}

a) $\mathrm{p}=3$

En caractéristique $3$, lorsque l'équation est mise sou\mbox{\begingroup\edef\encoutersep{\the\dimexpr.5\pgflinewidth\relax}\setbox0=\hbox{s}\setbox2=\hbox{\deleted{r}}\edef\encwda{\the\wd0}\edef\enchta{\the\ht0}\edef\encdpa{\the\dp0}\edef\encwdb{\the\wd2}\edef\enchtb{\the\ht2}\edef\encdpb{\the\dp2}\rlap{\smash{\tikz[baseline=0pt]{\begin{scope}[shift={(\the\dimexpr \encwda/2\relax,\the\dimexpr -\encdpa/2+\enchta/2\relax)}]\coordinate (enc-ending-center) at (0pt,0pt);\handwritten{\draw[overlay,thin,gray] ([xshift=-.1em,yshift=.85em]enc-ending-center) .. controls +(-.35em,.32em) and +(-.45em,.82em) .. ([xshift=-.67em,yshift=-.05em]enc-ending-center) .. controls +(.2em,-1em) and +(.8em,-.48em) .. ([xshift=.59em,yshift=-.1em]enc-ending-center) .. controls +(0,.75em) and +(.42em,.34em) .. ([xshift=-.1em,yshift=.55em]enc-ending-center);}\end{scope}\pgfresetboundingbox\path[use as bounding box] (0pt,\the\dimexpr-\encdpa\relax) rectangle (\encwda,\enchta);}}}\rlap{\kern\the\dimexpr \encwda/2-\encwdb/2\relax\raisebox{0pt}[0pt][0pt]{\copy2}}\copy0\endgroup} forme
% archive sheet number: 23
